\section{Interpolation Error}\label{sec:InterpolationError}
We recall perspective-correct interpolation (\cref{sec:PerspectiveCorrect}), derive the error of affine interpolation along an edge (\cref{sec:EdgeError}), and bound the error after subdivision (\cref{sec:Subdivision}).

\subsection{Perspective-Correct Interpolation}\label{sec:PerspectiveCorrect}
A triangle has the vertices $\pnt{p}_0$, $\pnt{p}_1$, $\pnt{p}_2$ in camera space with the texture coordinates $\vect{u}_0$, $\vect{u}_1$, $\vect{u}_2$ in texels.
The projection maps the vertex $\pnt{p}_k$ to the homogeneous coordinates $[x_k, y_k, z_k, w_k]^\top$ and to the screen position $[x_k/w_k, y_k/w_k]^\top$.
For a perspective projection, $w_k$ is the depth of $\pnt{p}_k$.
The rasterizer computes the screen-space barycentric coordinates $\lambda_0$, $\lambda_1$, $\lambda_2$ of a pixel with edge functions~\cite{Pineda88PAP}.
Affine interpolation weights the texture coordinates with them,
\begin{equation}
    \vect{u}_\mathrm{aff} = \sum_{k=0}^{2} \lambda_k \vect{u}_k.
\end{equation}
The texture coordinates are affine functions in camera space, however, not in screen space.
The correct texture coordinates $\vect{u} = \sum_{k} \mu_k \vect{u}_k$ use the weights~\cite{Heckbert91IPT}
\begin{equation}
    \mu_k = \frac{\lambda_k / w_k}{\sum_{j=0}^{2} \lambda_j / w_j}.
    \label{eq:Weights}
\end{equation}
\Cref{lst:Interpolation} implements both interpolations.

% Names match the symbols of the text: lambda for $\lambda_k$, w for $w_k$, u for $\vect{u}_k$, mu for $\mu_k$.
% mathescape typesets $...$ in the comments of the listing.
\begin{lstlisting}[language=C++, float=t, mathescape=true, morekeywords={float2, float3},
    label=lst:Interpolation, caption={\captiontitle{Texture Coordinates of a Pixel}
    Affine interpolation weights the texture coordinates with the barycentric coordinates $\lambda_k$,
    perspective-correct interpolation with the weights $\mu_k$ of \cref{eq:Weights}.}]
// Barycentric coordinates $\lambda_k$ as weights.
float2 affine(float3 lambda, const float2 u[3]) {
    return lambda.x * u[0] + lambda.y * u[1]
         + lambda.z * u[2];
}

// Weights $\mu_k$; w holds the depths $w_k$.
float2 perspective(float3 lambda, float3 w,
                   const float2 u[3]) {
    float3 q = lambda / w;
    float3 mu = q / (q.x + q.y + q.z);
    return affine(mu, u);
}
\end{lstlisting}


\subsection{Error Along an Edge}\label{sec:EdgeError}
We analyze the edge from $\pnt{p}_0$ to $\pnt{p}_1$ and order its endpoints such that $w_0 \le w_1$.
A pixel on the edge has the barycentric coordinates $\lambda_0 = 1 - t$, $\lambda_1 = t$, and $\lambda_2 = 0$, where $t\in[0,1]$ is its position on the edge in screen space.
The pixel shows the point $(1 - s)\,\pnt{p}_0 + s\,\pnt{p}_1$ with $s = \mu_1$.
With the \term{depth ratio} $r = w_1 / w_0 \ge 1$, \cref{eq:Weights} gives
\begin{equation}
    s = \frac{t / w_1}{(1 - t) / w_0 + t / w_1} = \frac{t}{r - (r - 1)\,t}.
\end{equation}
\Cref{fig:EdgeDiagram} illustrates this relation for $r = 4$.
Affine interpolation uses $t$ instead of $s$.
It is off by $\Vert \vect{u}_\mathrm{aff} - \vect{u} \Vert = (t - s)\,\ell$, where $\ell = \Vert \vect{u}_1 - \vect{u}_0 \Vert$ is the length of the edge in the texture.
The relative error
\begin{equation}
    \delta(t) = t - s = \frac{(r - 1)\,t\,(1 - t)}{r - (r - 1)\,t}
\end{equation}
vanishes at both endpoints.
It peaks where its derivative vanishes, at
\begin{equation}
    t^\ast = \frac{\sqrt{r}}{\sqrt{r} + 1},
    \qquad
    \delta_\mathrm{max} = \delta(t^\ast) = \frac{\sqrt{r} - 1}{\sqrt{r} + 1}.
    \label{eq:MaxError}
\end{equation}
The \SupplementalName{} gives the derivation.
The maximum depends only on the depth ratio: neither on the size of the triangle on the screen nor on the field of view.
For $r = 4$, the error peaks at two thirds of the edge on the screen and amounts to a third of the edge.
In our experiments, the error inside a triangle never exceeds the largest error along its edges (see the \SupplementalName).

\begin{figure}[t]
    \centering
    % Side view of an edge on the floor of the teaser (Sec. 3.2): the camera sits at the origin,
% the depth grows to the right. TikZ computes the exact geometry from the formulas of the paper.
\begin{tikzpicture}[x=16mm, y=22mm, >=stealth]
    \pgfmathsetmacro{\ratio}{4}      % depth ratio r with the depths w_0 = 1 and w_1 = r
    \pgfmathsetmacro{\screen}{0.8}   % depth of the screen
    \draw[->, PaletteGray] (0, 0) -- (4.35, 0) node[above left, black] {depth};
    \draw[PaletteGray, line width=1pt] (\screen, -0.95) -- (\screen, 0.12) node[above, black] {screen};
    \draw[line width=1pt] (1, -1) -- (\ratio, -1);
    % Rays through equal steps of t on the screen (black) hit the edge at s (blue);
    % affine interpolation assumes s = t (orange).
    \foreach \t in {0, 0.25, 0.5, 0.75, 1} {
        \pgfmathsetmacro{\s}{\t / (\ratio - (\ratio - 1) * \t)}
        \draw[PaletteGray] (0, 0) -- ({1 + (\ratio - 1) * \s}, -1);
        \fill (\screen, {-\screen * (1 - (1 - 1 / \ratio) * \t)}) circle (1pt);
        \fill[PaletteOrange] ({1 + (\ratio - 1) * \t}, -1) circle (1.6pt);
        \fill[PaletteBlue] ({1 + (\ratio - 1) * \s}, -1) circle (1.6pt);
    }
    \fill (0, 0) circle (1.5pt) node[above] {camera};
    \node[anchor=north east] at (\screen - 0.02, -0.82) {$t = 0$};
    \node[anchor=south west] at (\screen + 0.02, -0.17) {$t = 1$};
    \node[below] at (1, -1) {$\pnt{p}_0$};
    \node[below] at (\ratio, -1) {$\pnt{p}_1$};
    % The error at t = 1/2.
    \pgfmathsetmacro{\shalf}{0.5 / (\ratio - (\ratio - 1) * 0.5)}
    \draw[decorate, decoration={brace, mirror, amplitude=3pt}]
        ({1 + (\ratio - 1) * \shalf}, -1.06) -- ({1 + (\ratio - 1) * 0.5}, -1.06)
        node[midway, below=3pt] {$\delta(\tfrac{1}{2})$};
\end{tikzpicture}

    \caption{\captiontitle{Side View of an Edge}
        Rays through equal steps of $t$ on the screen (black) hit the edge at unequal steps of $s$ (blue).
        Affine interpolation assumes equal steps on the edge (orange).
        The brace marks the error $\delta(\tfrac{1}{2}) = 0.3$ for the depth ratio $r = 4$.
    }
    \label{fig:EdgeDiagram}
\end{figure}

\subsection{Subdivision}\label{sec:Subdivision}
Subdivision shortens the edges and reduces their depth ratios.
We split an edge into $n$ pieces of equal length in camera space.
The piece next to the camera has the largest depth ratio, $1 + (r - 1)/n$, and the length $\ell / n$ in the texture.
Since $(\sqrt{r} - 1)/(\sqrt{r} + 1) = (r - 1)/(\sqrt{r} + 1)^2 \le (r - 1)/4$, \cref{eq:MaxError} bounds the largest error $E_n$ of all pieces by
\begin{equation}
    E_n \le \frac{\ell}{n} \cdot \frac{r - 1}{4n} = \frac{\ell\,(r - 1)}{4n^2}.
    \label{eq:SubdivisionBound}
\end{equation}
The error falls quadratically with $n$.
The bound becomes tight for fine subdivisions, where the depth ratios of the pieces approach one.
To keep the error below a threshold $\varepsilon$, we need
\begin{equation}
    n \ge \sqrt{\frac{\ell\,(r - 1)}{4\varepsilon}}
    \label{eq:Pieces}
\end{equation}
pieces.
The diagonal of the floor of \cref{fig:Teaser} has the length $\ell = 362$~texels and the depth ratio $r = 4$.
For half a texel, $\varepsilon = 1/2$, \cref{eq:Pieces} requires $n \ge 23.3$, i.e., $24\times24$ cells or \num{1152} triangles.
